% Generated by tex/build_swebok_structure.py from tex/chapters/ch01/03_ソフトウェア要求の基礎.tex
\subsubsection{ソフトウェア要求の定義}

\term{ソフトウェア要求}とは、利用者が問題を解決したり目的を達成したりするために必要な条件または能力である。また、契約、標準、仕様、法令などを満たすために、システムや構成要素が備えるべき条件または能力でもある。さらに、それらを文書やモデルとして表したものも要求と呼ぶ。

より実務的に言えば、ソフトウェア要求とは「何を満たせば、現実世界の問題を解決したと言えるか」を表す合意である。ここでいう「問題」は、業務の自動化、既存ソフトウェアの欠点の修正、機器の制御、利用者体験の改善、規制への対応など、さまざまである。

\begin{definitionbox}ソフトウェア要求とは、ソフトウェアまたはソフトウェアを作るプロジェクトが満たすべき条件・能力・制約を、関係者が共有できるように表したものである。\end{definitionbox}

要求は通常、顧客、利用者、業務担当者から出る。しかし、それだけではない。規制当局、標準、既存システム、運用部門、サポート部門、開発組織、プロジェクト計画そのものも要求の源泉になる。
